%%=============================================================================
%% Inleiding
%%=============================================================================

\chapter{\IfLanguageName{dutch}{Inleiding}{Introduction}}%
\label{ch:inleiding}

In de hedendaagse logistiek is de opkomst van autonome mobiele robots (AMR's) niet meer weg te denken. Voor de ontwikkeling van deze robots wordt steeds vaker gebruikgemaakt van \textit{digital twins} in simulatieplatformen zoals NVIDIA Isaac Sim. Dit stelt ontwikkelaars in staat om algoritmen te testen in een veilige, virtuele omgeving voordat ze op dure hardware worden ingezet. Echter, de overgang van de virtuele naar de fysieke wereld — de \textit{sim-to-real gap} — blijft een kritiek punt. Waar onderzoek zich vaak richt op visuele of fysische verschillen (zoals lichtinval of frictie), blijkt de technische softwarematige integratie tussen de simulator en de robotica-middleware (ROS 2) een fundamentelere barrière te vormen.

Dit onderzoek richt zich op de integratie van NVIDIA Isaac Sim met de ROS 2 Nav2-stack voor een autonome magazijnrobot. Gedurende het traject verschoof de focus van een functionele prestatievergelijking naar een diepgaande analyse van de technische barrières die optreden bij het synchroniseren van coördinatenstelsels tussen deze platformen.

\section{\IfLanguageName{dutch}{Probleemstelling}{Problem Statement}}%
\label{sec:probleemstelling}

De overgang van een virtueel model naar een fysiek werkende robot wordt gehinderd door fundamentele verschillen in hoe simulators en robotica-frameworks data interpreteren. Tijdens de uitvoering van dit onderzoek bleek dat de virtuele robot, ondanks een visueel correcte opzet van de digitale tweeling, niet in staat was om via de SLAM Toolbox een ruimte in kaart te brengen. Dit resulteerde in een zogenaamd "rubberbanding"-effect, waarbij de robot herhaaldelijk naar de startpositie teleporteerde zodra sensordata werd verwerkt. De hoofdoorzaak van deze problematiek ligt in een mismatch tussen het coördinatenstelsel van Isaac Sim (Omniverse) en de gestandaardiseerde referentiestelsels van ROS 2 (REP-103 en REP-105).

De doelgroep voor dit onderzoek bestaat uit R\&D-ingenieurs en robotica-ontwikkelaars werkzaam in de logistieke automatisering die overwegen NVIDIA Isaac Sim te implementeren in hun ontwikkelcyclus. Specifiek biedt dit werk cruciale inzichten voor ontwikkelaars die gebruikmaken van de \texttt{isaac\_ros\_bridge} en tegen integratieproblemen aanlopen bij de configuratie van transformatiebomen ($TF$-trees).

\section{\IfLanguageName{dutch}{Onderzoeksvraag}{Research question}}%
\label{sec:onderzoeksvraag}

De centrale onderzoeksvraag van deze bachelorproef is als volgt geformuleerd:

\begin{quote}
    \textit{In welke mate vormen de technische discrepanties in coördinatenstelsels tussen NVIDIA Isaac Sim en ROS 2 een barrière voor de realisatie van een functionele digital twin voor autonome navigatie?}
\end{quote}

Om deze hoofdvraag te beantwoorden, worden de volgende deelvragen behandeld:
\begin{itemize}
    \item Wat zijn de theoretische vereisten voor een stabiele sim-to-real transitie bij AMR's in een magazijncontext?
    \item Welke invloed hebben afwijkingen in coördinatentransformaties ($TF$-frames) op de nauwkeurigheid van SLAM-algoritmen in een gesimuleerde omgeving?
    \item Waarom slaagt de Isaac ROS Bridge er in de huidige configuratie niet in om de odometrie-data consistent te synchroniseren met de ROS 2 Nav2-stack?
\end{itemize}

\section{\IfLanguageName{dutch}{Onderzoeksdoelstelling}{Research objective}}%
\label{sec:onderzoeksdoelstelling}

Het beoogde resultaat van deze bachelorproef is tweeledig:
\begin{enumerate}
    \item \textbf{Een fysiek prototype:} De realisatie van een autonome robot die in de fysieke wereld succesvol kan navigeren en mappen met ROS 2.
    \item \textbf{Een technisch analyserapport:} Een gedetailleerde documentatie van de geconstateerde integratiefouten in de virtuele omgeving.
\end{enumerate}

Het succes van deze proef wordt niet uitsluitend afgemeten aan een perfect werkende simulatie, maar aan de identificatie van de root-cause van de systeemfout. Het doel is om door middel van een foutenanalyse concrete aanbevelingen te formuleren voor de robotica-community om de softwarematige kloof tussen Omniverse-simulaties en ROS 2 te overbruggen.

\section{\IfLanguageName{dutch}{Opzet van deze bachelorproef}{Structure of this bachelor thesis}}%
\label{sec:opzet-bachelorproef}

% Het is gebruikelijk aan het einde van de inleiding een overzicht te
% geven van de opbouw van de rest van de tekst. Deze sectie bevat al een aanzet
% die je kan aanvullen/aanpassen in functie van je eigen tekst.

De rest van deze bachelorproef is als volgt opgebouwd:

In Hoofdstuk~\ref{ch:stand-van-zaken} wordt een overzicht gegeven van de stand van zaken binnen het onderzoeksdomein, waarbij dieper wordt ingegaan op AMR-technologieën en de ROS 2-standaarden REP-103 en REP-105.

In Hoofdstuk~\ref{ch:methodologie} wordt de methodologie toegelicht. Hierin wordt de opbouw van zowel de fysieke robot als de digitale tweeling besproken, inclusief de gebruikte diagnostische technieken voor de $TF$-tree analyse.

In Hoofdstuk~\ref{ch:resultaten} (Resultaten en Discussie) worden de bevindingen van de fysieke testen gepresenteerd en wordt een diepgaande analyse gegeven van het falen van de SLAM-fase in de simulator.

In Hoofdstuk~\ref{ch:conclusie}, tenslotte, wordt de conclusie gegeven en een antwoord geformuleerd op de onderzoeksvragen. Daarbij worden ook aanbevelingen gedaan voor toekomstige integraties tussen NVIDIA Isaac Sim en ROS 2.